\section{模型评价与改进}
\label{sec:evaluation}
\subsection{统一模型与分层求解的特点}
\subsubsection{从单架次属性到系统执行的连续论证}
全文以同一节点、货箱和装备数据为基础，用统一几何、逐段载荷、交接完成和能源恢复计算生成任务属性。问题一在这些属性上决定精确组批；问题二增加有限飞机和电池；问题三把真实阶段轨迹转为通信需求；问题四冻结完整日程，决定组织分组和资源配置。每一问都有明确新增的决策对象，避免把四个答案写成互不兼容的模型。

这种组织首先保证结果可衔接。一个箱号始终对应同一质量、体积、目的地与期限，任务在表格、甘特图和验证记录中保持同一标识；能源变化通过充电长度传到后继任务，通信变化通过提供方窗口传到运输时刻，分组则通过禁止共享改变配置峰值。系统层面的代价因此能回溯到具体物理或资源关系。

\subsubsection{按问题结构选择计算方法}
安全组批具有服务区可分离和计数状态结构，精确模型与动态规划既给出方案，也支持最少架次证明。大规模运输调度以高质量可行搜索产生完整日程，再用独立列松弛提供质量区间。通信部署涉及非线性几何和离散服务关系，采用候选筛查、事件区间、联合整数规划与固定结构线性精修。固定日程配置则利用区间结构与少量合法分区精确完成。

\begin{table}[htbp]\centering\small
\caption{各问题求解结构与主要证据的对应}\label{tab:evidence_levels}
\begin{tabularx}{\linewidth}{p{1.5cm}LLL}\toprule
问题 & 可利用的结构 & 方案获得方式 & 主要证据\\\midrule
单点组批 & 分区可分离、等价货箱计数 & 精确覆盖与状态递推 & 架次下界闭合、独立目标一致\\
运输调度 & 任务属性、双资源接续 & 邻域重构与事件排程 & 全箱重放、全局有效上下界\\
连续通信 & 状态区间、固定结构线性时序 & 候选部署与联合时序优化 & 连续区间/端点、固定结构LP\\
独立配置 & 半开区间、不可拆关联 & 峰值构造与完整分区枚举 & 资源链、匹配、向量支配\\\bottomrule
\end{tabularx}\end{table}

模型复杂程度与问题规模相匹配，是本题获得可执行解和可核验结论的重要条件。固定日程的最少配置不需要随机优化；连续位置的联合问题也不因最后一层是线性规划，就被等同为全局线性问题。方法分工使计算资源集中在真正的组合与几何难点。

\subsection{从数值结果提炼救援决策规律}
\subsubsection{载荷、能源和时间的相互作用}
混合机型在相同18架次下比仅C型方案节能约20.02\%，说明最大载荷不是唯一选择依据。进入有限资源日程后，飞机与电池不同步使工期取决于关键接续；八架运输机95.76\%的作业占用率，也使简单换序的改善空间受限。通信加入之后，位置和服务交接又决定关键运输任务能够多早启动。

因此，三类常见简化分别遗漏不同机制：仅按额定载荷会忽略山区安全能力，仅按距离会遗漏交接与能源恢复，仅按静态通信覆盖会遗漏中继到位与提供方切换。本文用具体任务、接续与区间证据逐层补足这些关系，从而获得同时满足物资、装备和通信要求的方案。

\subsubsection{共享资源与独立组织的量化价值}
相同日程在全局共享下配置27件，两组独立28件，三组独立34件。增加来自跨组错峰复用的丧失，不来自新增运输需求。所选两组需求还逐分量支配其他合法分区，使库存变化下的选择稳定性可以由向量单调性证明。资源件数、型号、工作量和组数同时展示，能直接服务于独立救援单元的装备清单设计。

\subsection{保障与恢复的实施路径}
\subsubsection{按变化来源选择最小调整层次}
安全余量变化先检查原箱组是否仍在可行域；充电和交接延迟先沿固定资源顺序修复；能量约束失效则重新组批或改型；传播损耗影响链路门限时使用相应保障备选，中继上线变化优先重新协调服务时段。这样，扰动不是仅作用于结果工期的系数，而是沿其实际影响的物理和调度层传播。

\begin{table}[htbp]\centering\small
\caption{模型结果对应的分级执行调整}\label{tab:execution_policy}
\begin{tabularx}{\linewidth}{>{\raggedright\arraybackslash}p{2.2cm}LL}\toprule
变化来源 & 优先检查量 & 对应措施\\\midrule
安全要求提高 & 原任务SOC及安全载荷边界 & 范围内沿用，越过临界点重新组批\\
充电与交接变慢 & 接续余量、后继最早开始与箱级时限 & 固定结构时刻修复\\
任务能耗增加 & 返航安全预算与恢复时间 & 先检查物理可行，再选择结构调整备选\\
传播条件变差 & 双跳门限、覆盖区间及切换端点 & 使用匹配保障方案并完整核验\\
中继上线延后 & 后继提供方就绪与关键交接约束 & 服务窗口和运输时刻联合修复\\
组织或库存变化 & 逐组需求向量与正缺口 & 重新核算配置并按类型补充\\\bottomrule
\end{tabularx}\end{table}

\subsubsection{与更广泛不确定优化的衔接}
本文的压力情景使用明确的参数倍率和单项延迟，能够定位临界约束与恢复代价。若未来积累可信的天气、故障与时延样本，可以进一步构造情景集合、分布鲁棒或滚动时域模型。相关灾后物流研究已探索需求分布信息不足、通信不确定和动态天气条件下的规划\cite{xin2025,dukkanci2026,li2025hurricane,zhang2026communication}；这些研究说明了扩展方向，但不被冒充为本题已经运行的算法或实地数据。

扩展时仍应保留现有接口：原始需求保持唯一，物理参数在任务层更新，变化后的任务回到资源与通信验证，最终再核算独立配置。当前精确范围为给定确定性物理规则及相应保存方案；任务来源、完整复算入口和不同层次证明范围集中列于附录。这样的组织既突出已完成工作的结果，也为引入新的实际条件留下清晰的计算位置。
